\documentclass{article}
\usepackage{amsmath, amsthm, amssymb}

\newtheorem{theorem}{Theorem}
\newtheorem{lemma}{Lemma}
\newtheorem{corollary}{Corollary}
\newtheorem{assumption}{Assumption}

\title{Mathematical Proofs of Stability for the CausalNerve Lyapunov Gate}
\author{CausalNerve Systems Theory Group}
\date{\today}

\begin{document}
\maketitle

\section{Introduction}
This document formalizes the stability bounds of the Online Causal Graph Revision (OCGR) module.

\section{Definitions and Assumptions}
Let $\mathcal{G}$ be the finite space of all directed graphs on $N$ nodes.
Let the energy function be $V: \mathcal{G} \to \mathbb{R}^+$.

\begin{assumption}[Boundedness]
There exists a finite $V_{max}$ such that $V(G) \le V_{max}$ for all $G \in \mathcal{G}$ generated by the system, and $V(G) \ge 0$.
\end{assumption}

\begin{assumption}[Strict Monotonicity]
The gating policy enforces that an edit sequence $(G_0, G_1, \dots, G_k)$ is valid if and only if $V(G_{i}) - V(G_{i+1}) \ge \epsilon$ for all $i$, where $\epsilon > 0$ is the strict margin.
\end{assumption}

\section{Theorems and Proofs}

\begin{theorem}[Finite Oscillation Bound]
Under Assumptions 1 and 2, any continuous sequence of graph edits is strictly bounded in length by $K_{max} \le \frac{V(G_0)}{\epsilon}$.
\end{theorem}

\begin{proof}
Let $(G_0, G_1, \dots, G_k)$ be a valid sequence of accepted edits. 
By Assumption 2, for each step $i$, $V(G_i) - V(G_{i+1}) \ge \epsilon$.
Summing over the sequence from $i=0$ to $k-1$:
\begin{align*}
\sum_{i=0}^{k-1} [V(G_i) - V(G_{i+1})] &\ge \sum_{i=0}^{k-1} \epsilon \\
V(G_0) - V(G_k) &\ge k \cdot \epsilon
\end{align*}
Since $V(G_k) \ge 0$ (Assumption 1), we have:
\begin{align*}
V(G_0) &\ge k \cdot \epsilon \\
k &\le \frac{V(G_0)}{\epsilon}
\end{align*}
Because $k$ is bounded by a finite constant, the number of sequential structural edits is finite. Thus, infinite oscillation (alarm storms resulting in continuous graph flipping) is mathematically impossible.
\end{proof}

\begin{theorem}[Asymptotic Stability in the Edit Graph]
Let the transition system be a directed graph $\mathcal{H}$ where nodes are causal graphs $G \in \mathcal{G}$, and a directed edge exists from $G_a \to G_b$ iff $o(G_a) = G_b$ and $V(G_b) < V(G_a) - \epsilon$. The system is asymptotically stable and converges to a sink node (equilibrium graph).
\end{theorem}

\begin{proof}[Proof Sketch]
The relation $G_a \to G_b$ implies a strict total ordering on the energy $V$. Therefore, the transition graph $\mathcal{H}$ is a Directed Acyclic Graph (DAG). 
Any random walk on a finite DAG must terminate at a sink node (a node with out-degree zero). Thus, the OCGR algorithm must halt at a local minimum graph $G_{eq}$ where no single edit operator $o \in \mathcal{O}$ can satisfy the Lyapunov gate. The state is asymptotically stable.
\end{proof}

\section{Open Problems and Limitations}

\subsection{Convergence Under Stochastic Alarms (Incomplete)}
The current proof relies on the assumption that $L_{leak}(G_t)$ is evaluated deterministically. In reality, $L_{leak}$ is computed over a sliding window of stochastic sensor data. 
\textbf{Limitation:} If the data window shifts, the landscape of $V(G)$ itself changes dynamically over time ($V_t(G)$). 
If the rate of change of the landscape $\frac{\partial V}{\partial t}$ is larger than the descent rate $\epsilon$, Theorem 1 breaks down, and the system can enter a slow drift oscillation. A formal proof of bounds under stochastic, time-varying landscapes requires martingale theory and remains an open problem.

\end{document}
